\section{Problemas y desafíos de SDS}

Aunque SDS ha logrado resultados impresionantes, también presenta algunos problemas inherentes.

\subsection{El problema Janus (problema multifacético)}

El modo de falla más famoso de SDS es el \textbf{Janus Problem} (nombrado por el dios romano bifronte Janus): los modelos 3D generados parecen ver hacia adelante desde cada perspectiva, lo que resulta en un solo objeto mostrando múltiples caras.

\begin{warningbox}{Causa raíz del problema Janus}
La causa fundamental del problema Janus reside en la contradicción entre el prior 2D y la consistencia 3D:

\begin{enumerate}
    \item Los modelos de difusión de imágenes se entrenan en imágenes de \textbf{una sola perspectiva}, por lo que no poseen el concepto de consistencia 3D.
    \item SDS optimiza los resultados de renderizado de cada perspectiva de forma independiente; cada perspectiva tiende a generar la imagen más ``típica'' de esa descripción textual.
    \item Para objetos con una cara claramente definida, como rostros humanos, la imagen más ``típica'' es casi siempre desde la perspectiva frontal.
    \item El resultado es que cada cara del modelo 3D se convierte en una cara frontal, generando múltiples rostros.
\end{enumerate}

En esencia, esto se debe a que los modelos de difusión 2D solo poseen un prior de \textbf{observación local} (imágenes de una sola perspectiva), mientras que la generación 3D requiere restricciones de \textbf{consistencia global}.
\end{warningbox}

\subsection{Compromiso entre el peso CFG y la diversidad}

Al utilizar SDS, generalmente es necesario establecer un peso de Classifier-Free Guidance (CFG) muy alto para obtener una salida limpia.

\begin{knowledgebox}{Efecto de doble filo de un alto peso de CFG}
Un peso de CFG elevado contrae la distribución generativa hacia una región estrecha que coincide altamente con la descripción textual:

\begin{itemize}
    \item \textbf{Ventaja}: mayor calidad de salida y mayor coherencia con la descripción textual.
    \item \textbf{Desventaja}: la diversidad disminuye drásticamente; todos los resultados generados tienden a parecerse entre sí.
\end{itemize}

Cómo mantener la diversidad al mismo tiempo que se garantiza la calidad constituye una pregunta abierta importante para SDS.
\end{knowledgebox}

\subsection{Otras limitaciones}

\begin{itemize}
    \item \textbf{Problema de sobreexposición}: las texturas generadas por SDS suelen ser excesivamente saturadas y carecen de naturalidad.
    \item \textbf{Optimización inestable}: debido a la perspectiva aleatoria y el ruido aleatorio en cada iteración, el proceso de optimización puede oscilar.
    \item \textbf{Costo computacional}: cada iteración requiere una propagación hacia adelante del modelo de difusión; todo el proceso de optimización necesita miles de iteraciones.
    \item \textbf{Calidad geométrica}: la geometría 3D generada a menudo no es suficientemente detallada, especialmente en las zonas ocultas.
\end{itemize}

\begin{warningbox}{De un prior local a una consistencia global: dificultad fundamental}
Todos los problemas de SDS pueden rastrearse a una dificultad fundamental única: utilizar un \textbf{prior de observación 2D local} para restringir la \textbf{estructura 3D global}. Este es un problema mal planteado (ill-posed): infinitas estructuras 3D pueden producir renderizados 2D que parezcan razonables a simple vista. Un solo prior 2D no es suficiente para eliminar esta ambigüedad; se requiere introducir conocimiento previo específico del dominio 3D.
\end{warningbox}

\subsection{Resumen de la sección}

Los principales desafíos de SDS incluyen el problema Janus, texturas sobreexpuestas y el ajuste del peso CFG. La raíz de estos problemas radica en que la localidad del prior 2D no puede restringir completamente la consistencia global de la estructura 3D. Trabajos posteriores (como las diversas versiones mejoradas de SDS: VSD, ISM, etc.) se dedican a mitigar estos problemas.

\newpage

